\section{Data and Partitioning}

\subsection{GLODAP v2.2023 dataset}
Observational data were extracted from the Global Ocean Data Analysis Project version 2 (GLODAP v2.2023) product \citep{Lauvset2023}. GLODAP provides internally consistent, quality-controlled measurements of ocean physical and biogeochemical properties from bottle samples collected on international research cruises. We extracted salinity ($S$), in situ temperature ($T$, $^\circ\mathrm{C}$), apparent oxygen utilisation ($\aou$, $\mu\mathrm{mol}\,\mathrm{kg}^{-1}$), total dissolved inorganic carbon ($\tco$, $\mu\mathrm{mol}\,\mathrm{kg}^{-1}$), as well as metadata including latitude, longitude, depth (m), sampling year, GLODAP region code, and cruise identifier.

\subsection{Quality control filtering}
Strict quality control filters were applied to ensure dataset integrity. Observations were retained only if all primary variables ($S$, $T$, $\aou$, $\tco$) and metadata fields were non-null, finite, and satisfied physically realistic physical boundaries:
\begin{align}
25.0 &< S < 42.0\,\mathrm{PSU}, \\
-2.5 &< T < 35.0\,^\circ\mathrm{C}, \\
1700.0 &< \tco < 2600.0\,\mu\mathrm{mol}\,\mathrm{kg}^{-1}, \\
\aou &> -50.0\,\mu\mathrm{mol}\,\mathrm{kg}^{-1}.
\end{align}
Following these QC filters, a total of 472,530 high-quality global observations were retained for downstream model fitting, validation, and evaluation.

\subsection{Basin definitions}
To account for distinct oceanographic circulation regimes, ventilation pathways, and water-mass characteristics \citep{Tomczak2002}, observations were partitioned into four mutually exclusive geographic basins based on latitude and GLODAP region definitions:
\begin{enumerate}[label=(\roman*)]
\item \textbf{Atlantic Ocean}: Region code 1, restricted to latitude $\ge -35.0^\circ$ ($N = 147{,}448$);
\item \textbf{Indian Ocean}: Region code 16, restricted to latitude $\ge -35.0^\circ$ ($N = 40{,}391$);
\item \textbf{Pacific Ocean}: Region code 8, restricted to latitude $\ge -35.0^\circ$ ($N = 187{,}443$); and
\item \textbf{Southern Ocean}: All observations with latitude $< -35.0^\circ$ across all longitudes ($N = 97{,}248$).
\end{enumerate}

\subsection{Predictor scaling and notation}
In symbolic regression and polynomial modeling, predictor scaling is critical for numerical stability. We scale $\aou$ by a factor of 100 ($\aou_{100} = \aou / 100$), ensuring that all three predictors ($S$, $T$, $\aou_{100}$) share comparable numerical magnitudes order-of-magnitude wise ($S \sim 35$, $T \sim 0$--$30$, $\aou_{100} \sim -0.5$--$4.0$). All candidate models operate directly on these physical variables without standard z-score normalization to preserve direct scientific interpretability.

\subsection{Temporal partitioning and external holdout}
To evaluate model stability and prevent temporal data leakage, data were partitioned chronologically into three disjoint subsets:
\begin{enumerate}[label=(\roman*)]
\item \textbf{Training set} (Years $< 2015$): $N = 378{,}507$ observations used exclusively for model coefficient fitting and symbolic expression generation;
\item \textbf{Temporal validation set} (Years $2015$--$2017$): $N = 53{,}474$ observations used for hyperparameter tuning, model family comparison, and Pareto candidate selection; and
\item \textbf{Locked external holdout} (Years $\ge 2018$): $N = 40{,}549$ observations strictly reserved for final temporal out-of-sample evaluation.
\end{enumerate}

Importantly, the post-2018 external holdout remained completely locked throughout candidate search and model selection. No hyperparameter adjustment, equation choice, or narrative optimization was performed using post-2018 data. Table~\ref{tab:sample_sizes} summarizes the sample size distributions across basins and temporal splits, and Table~\ref{tab:predictor_distributions} details predictor summary statistics.

\begin{table}[htbp]
\centering
\caption{Quality-controlled GLODAP v2.2023 sample sizes after basin assignment. Training years are $<2015$; temporal validation years are $2015$--$2017$; the locked external holdout is $\mathrm{year}\ge 2018$.}
\label{tab:sample_sizes}
\begin{tabular}{lrrrr}
\toprule
Basin & $N$ & Training & Temporal validation & External holdout \\
\midrule
Atlantic & 147{,}448 & 125{,}762 & 11{,}436 & 10{,}250 \\
Indian & 40{,}391 & 32{,}498 & 3{,}764 & 4{,}129 \\
Pacific & 187{,}443 & 138{,}202 & 31{,}791 & 17{,}450 \\
Southern Ocean & 97{,}248 & 82{,}045 & 6{,}483 & 8{,}720 \\
Total & 472{,}530 & 378{,}507 & 53{,}474 & 40{,}549 \\
\bottomrule
\end{tabular}
\end{table}

\begin{table}
\caption{Predictor and target distributions by basin.}
\label{tab:predictor_distributions}
\resizebox{\textwidth}{!}{%
\begin{tabular}{lrrrrrrr}
\toprule
Basin & Variable & N & Mean & SD & Q25 & Median & Q75 \\
\midrule
Unclassified & salinity & 512384 & 34.665 & 0.899 & 34.483 & 34.694 & 34.926 \\
Unclassified & temperature & 512384 & 6.707 & 7.430 & 1.626 & 3.562 & 10.121 \\
Unclassified & aou & 512384 & 101.873 & 85.448 & 29.322 & 85.784 & 163.690 \\
Unclassified & tco2 & 512384 & 2187.727 & 104.867 & 2127.000 & 2182.000 & 2261.300 \\
Unclassified & depth & 512384 & 1213.613 & 1380.175 & 139.000 & 613.000 & 1943.925 \\
\bottomrule
\end{tabular}%
}
\end{table}

